% ===== 原文 pp.34–36（PDF pp.35–37）=====

% [原文 p.34]
给定的 $f:U(1)^{N-k}\to T_F$ 并不能唯一地确定整数 $f_i^\alpha$——存在把 $f_i^\alpha$ 平移 $\sum_a Q_i^a m_a^\alpha$（$m_a^\alpha$ 为某些整数）的歧义。但得益于真空方程 $\sum_{i=1}^{N}Q_i^a|\phi_i|^2=\zeta^a$，这只会把该 $\mu_F^\alpha(\phi)$ 平移一个常数 $\sum_a\zeta^a m_a^\alpha$。

$\mu_F$ 在 $X_\zeta$ 上的像 [7] 是 $\ii\ft_F^*$ 中的一个凸多面体，我们称其为 $X_\zeta$ 的\emph{矩多面体} $P_\zeta$。$P_\zeta$ 的内点对应于所有分量 $\phi_i$ 都非零的 $[\phi]\in X_\zeta$，而边界点则对应于某些 $\phi_i$ 为零的那些点。
% [Tutor Note: 原文此处 "the isomorophism"（p.34 第 8 行）是 "isomorphism" 的笔误；另外 p.34 末段
% "either of the the orbifold phase" 多打了一个 "the"（见印刷 p.34 第 20 行）。译文按正确含义处理。]

例如，考虑 $\CP^{N-1}$ 模型的 $\zeta>0$ 相。作为同构 $f:U(1)^{N-1}\to T_F=U(1)^N/G$，我们可以取 $(c_1,\ldots,c_{N-1})\mapsto[(c_1,\ldots,c_{N-1},1)]$。此时矩映射为 $\mu_F(\phi)=(|\phi_1|^2,\ldots,|\phi_{N-1}|^2)$，其中 $\phi$ 满足 $\sum_{i=1}^{N}|\phi_i|^2=\zeta$。由于每个 $|\phi_i|^2$ 都非负，我们看到 $(\eta^1,\ldots,\eta^{N-1})\in\ii\ft_F^*\cong\R^{N-1}$ 属于 $P_\zeta$ 当且仅当 $\eta^\alpha\ge 0$ 且 $\eta^1+\cdots+\eta^{N-1}\le\zeta$。也就是说，$P_\zeta$ 是 $\R^{N-1}$ 中标准的 $(N-1)$-单纯形按 $\zeta$ 缩放后的结果。$N=2$ 与 $N=3$ 的情形见图 6，图中把在边界超平面上消失的场标注了出来。

\begin{figure}[htbp]
\centering
\includegraphics[width=0.9\textwidth]{fig6.png}
\figtitle{6}{$\CP^1$（左）与 $\CP^2$（右）的矩多面体。}
\end{figure}

我们再来描述其他例子中的矩多面体。对于 $\mathrm{T}^{U(1)}_{N,d}(0)$，
\begin{align}
    &f:(c_0,c_1,\ldots,c_{N-1})\in U(1)^N\longmapsto[(c_0,c_1,\ldots,c_{N-1},1)]\in U(1)^{N+1}/G,
    \notag\\
    &P_\zeta=\Big\{(\eta^0,\eta^1,\ldots,\eta^{N-1})\ \Big|\ \eta^\alpha\ge 0,\ 
    \eta^1+\cdots+\eta^{N-1}-d\eta^0\le\zeta\Big\}, \tag{4.43}
\end{align}
而对于 $\mathrm{H}_n$，
\begin{align}
    &f:(c_1,c_2)\in U(1)^2\longmapsto[(c_1,1,c_2,1)]\in U(1)^4/G, \notag\\
    &P_\zeta=\Big\{(\eta^1,\eta^2)\ \Big|\ \eta^1\ge 0,\ 0\le\eta^2\le\zeta^2,\ 
    \eta^1+n\eta^2\le\zeta^1+n\zeta^2\Big\}. \tag{4.44}
\end{align}
在两个例子中，$P_\zeta$ 的表达式在轨形相（对 $\mathrm{T}^{U(1)}_{N,d}$ 是 $\zeta<0$，对 $\mathrm{H}_n$ 是相 II）与几何相（对 $\mathrm{T}^{U(1)}_{N,d}$ 是 $\zeta>0$，对 $\mathrm{H}_n$ 是相 I）中都成立。见图 7 与图 8，那里我们同样把在边界超平面上消失的场标注了出来。
% [Tutor Note: (4.43)(4.44) 中的不等式都可逐条由 D-项方程与非负性 |\phi_i|^2\ge0 反推，是最好的
% 自查锚点。对 H_n：把 (4.34)（印刷 p.29）的 |\phi_1|^2+|\phi_2|^2-n|\phi_4|^2=\zeta^1、
% |\phi_3|^2+|\phi_4|^2=\zeta^2 与 \eta^1=|\phi_1|^2、\eta^2=|\phi_3|^2 联立，则 \eta^1\ge0（|\phi_1|^2\ge0）、
% 0\le\eta^2\le\zeta^2（|\phi_3|^2\ge0 与 |\phi_4|^2=\zeta^2-\eta^2\ge0）、
% \eta^1+n\eta^2\le\zeta^1+n\zeta^2（|\phi_2|^2=\zeta^1+n\zeta^2-\eta^1-n\eta^2\ge0）三条恰好一一对应。
% 对 T^{U(1)}_{N,d}(0)：由印刷 p.28 的 -d|p|^2+|x|^2=\zeta 与 \eta^0=|p|^2、\eta^\alpha=|x_\alpha|^2
% （\alpha=1,\ldots,N-1），并注意最后一个场 |x_N|^2=\zeta+d\eta^0-\sum_{\alpha=1}^{N-1}\eta^\alpha\ge0，
% 即得 \eta^1+\cdots+\eta^{N-1}-d\eta^0\le\zeta。]

% [原文 p.35]
\begin{figure}[htbp]
\centering
\includegraphics[width=0.9\textwidth]{fig7.png}
\figtitle{7}{$\C^2/\Z_d$（左）与 $\mathrm{tot}[\cO(-d)\to\CP^1]$（右）的矩多面体。}
\end{figure}

\begin{figure}[htbp]
\centering
\includegraphics[width=0.9\textwidth]{fig8.png}
\figtitle{8}{$\WP^2_{1,1,n}$（左）与 $\mathbb{F}_n$（右）的矩多面体。}
\end{figure}

对我们来说最重要的一点是，$\mu_F:X_\zeta\to P_\zeta$ 是一个拉格朗日环面纤维化。在 $P_\zeta$ 的内点之上，它是纤维为 $T_F$ 的光滑环面纤维化——事实上是一个主 $T_F$ 丛——因为所有 $\phi_i$ 都非零。在 $P_\zeta$ 的边界处，随着某些 $\phi_i$ 变为零，一部分 $T_F$ 纤维退化。这正是我们在 $\CP^1$ 情形（见图 1）已经见到的现象，而 GLSM 给出了它的系统推广。如同 $\CP^1$ 的情形一样，我们要问：对环面纤维作 T-对偶时会发生什么，来自奇异纤维的量子修正又是什么。这将是下一节的主题。

作为复流形，或作为代数簇，$X_\zeta$ 可以描述为 $V$ 按由 $\zeta$ 决定的稳定性条件对复化规范群 $G_{\C}$ 所作的几何不变量理论商。从集合论上说，它就是 $V$ 中不稳定轨迹 $F_\zeta$ 的补集对 $G_{\C}$ 的商：
\begin{equation}
    X_\zeta=(V\setminus F_\zeta)/G_{\C}. \tag{4.45}
\end{equation}
轨迹 $F_\zeta$ 是线性子空间的并
\begin{equation}
    F_\zeta=\bigcup_{I\in\cI_\zeta}\{\phi_i=0,\ \ \forall i\in I\} \tag{4.46}
\end{equation}
其中 $\cI_\zeta$ 是 $\{1,\ldots,N\}$ 的一族非空子集。对 $v\in\ii\fg$，令 $I_v\subset\{1,\ldots,N\}$%
% [原文 p.36]
为使 $Q_i(v)>0$ 的那些 $i$ 的集合。那么，
\begin{equation}
    \cI_\zeta:=\Big\{\varnothing\ne I\subset\{1,\ldots,N\}\ \Big|\ 
    I_v\subset I\ \text{对某个满足}\ \zeta(v)>0\ \text{的}\ v\in\ii\fg\Big\}. \tag{4.47}
\end{equation}
% [Tutor Note: (4.47) 是 Hilbert--Mumford 判据在 G=U(1)^k 情形的具体形式：I\in\cI_ζ 当且仅当存在
% v\in i\fg 使得 ζ(v)>0 且所有满足 Q_i(v)>0 的指标 i 都落在 I 中。直观理由：在真空流形上
% \langle\mu(\phi),v\rangle=\sum_i Q_i(v)|\phi_i|^2=\zeta(v)>0，故 ζ(v)>0 时至少有一个正权场非零，
% 把 I_v 整体置零的坐标子空间根本不会出现在 \mu^{-1}(ζ) 里，必须先挖掉；由 Kempf--Ness 定理，
% 这样挖掉之后的 GIT 商 (V\setminus F_ζ)/G_{\C} 与辛商 (4.40) 一致。注意 (4.47) 中的 v 取遍整个
% 李代数 i\fg（所有实权方向），而不只是 (4.50) 的那 N 个 v^i——这是它与 (4.54) 的差别所在。]

这一描述清楚地表明，环面味对称群 $T_F$ 的作用可以延拓为其复化 $T_{F\C}$ 的全纯作用。事实上，满足所有 $i$ 都有 $\phi_i\ne 0$ 的点（即 $P_\zeta$ 的内点在 $\mu_F$ 下的原像）构成 $T_{F\C}$ 的一条单一自由轨道，它在 $X_\zeta$ 中稠密开。具有这种环面作用的代数簇称为\emph{环面簇}。

环面簇的标准描述是用\emph{扇}（fan）的语言给出的。由提供给我们的数据，它如下给出。设 $\mathrm{M}$ 与 $\mathrm{N}$ 分别是 $T_F$ 的权格与余特征格：
\begin{align}
    \mathrm{M} &= \Hom(T_F,U(1)), \tag{4.48}\\
    \mathrm{N} &= \Ker\,(v\in\ii\ft_F\mapsto\exp(2\pi\ii v)\in T_F). \tag{4.49}
\end{align}
它们自然地互为对偶，且 $\mathrm{M}\otimes\R=\ii\ft_F^*$、$\mathrm{N}\otimes\R=\ii\ft_F$。我们定义 $v^1,\ldots,v^N\in\mathrm{N}$ 为
\begin{equation}
    v^i=[(0,\ldots,0,1,0,\ldots,0)]\qquad(\text{第 }i\text{ 个位置为 }1), \tag{4.50}
\end{equation}
使得 $\exp(\ii\sum_{i=1}^{N}\theta_i v^i)=[(\ee^{\ii\theta_1},\ldots,\ee^{\ii\theta_N})]$。设 $v_\alpha^i$ 是 $v^i$ 在由 (4.41) 确定的基下的分量，即 $v^i=f(v_1^i,\ldots,v_{N-k}^i)$。由于 $t\mapsto\exp(\ii tv^i)$ 的周期为 $2\pi$，分量 $v_\alpha^i$ 必须全为整数。事实上，它们给出了 (4.41) 的逆：
\begin{equation}
    f^{-1}:[(\omega_1,\ldots,\omega_N)]\in T_F\longmapsto
    \Big(\prod_i\omega_i^{v_1^i},\ldots,\prod_i\omega_i^{v_{N-k}^i}\Big)\in U(1)^{N-k}. \tag{4.51}
\end{equation}
由于它必须是良定义的映射，我们有
\begin{equation}
    \sum_{i=1}^{N}v_\alpha^i Q_i^a=0. \tag{4.52}
\end{equation}
由于 (4.41) 与 (4.51) 互为逆映射，我们有
\begin{equation}
    \sum_{i=1}^{N}v_\beta^i f_i^\alpha=\delta_\beta^\alpha,\qquad
    \sum_{\alpha=1}^{N-k}f_j^\alpha v_\alpha^i=\delta_j^i+\sum_{a=1}^{k}Q_j^a\xi_a^i
    \quad(\text{对某些}\ \xi_a^i). \tag{4.53}
\end{equation}
% [Tutor Note: (4.52) 正是 (4.51) 良定义的条件：若把代表元 \omega_i 换成 \ee^{\ii Q_i^a\varepsilon}\omega_i
% （同一个 T_F=U(1)^N/G 中的点），则 \prod_i\omega_i^{v_\alpha^i} 乘以
% \exp(\ii\varepsilon\sum_i v_\alpha^i Q_i^a)，要求它恒为 1 即得 (4.52)。同理 (4.53) 第一式即 f^{-1}\circ f=\id，
% 第二式即 f\circ f^{-1} 与恒等映射只差规范方向 \sum_a Q_j^a\xi_a^i。另外原文 "1 at the i-th intry" 中
% "intry" 是 "entry" 之误。]
注意，若 $\phi\in V$ 满足真空方程 $\sum_i Q_i^a|\phi_i|^2=\zeta^a$，则
\begin{align}
    \langle\mu_F(\phi),v^i\rangle
    &=\sum_\alpha\Big(\sum_j f_j^\alpha|\phi_j|^2\Big)v_\alpha^i \notag\\
    &=\sum_j\Big(\delta_j^i+\sum_a Q_j^a\xi_a^i\Big)|\phi_j|^2
    =|\phi_i|^2+\sum_a\zeta^a\xi_a^i. \tag{4.54}
\end{align}
% [Tutor Note: (4.54) 的两步都值得核对：第一步用 (4.42) 把 \mu_F^\alpha(\phi) 展开并缩并 v_\alpha^i；
% 第二步用 (4.53) 第二式替换 \sum_\alpha f_j^\alpha v_\alpha^i，最后一步把 \sum_j Q_j^a|\phi_j|^2 换成 \zeta^a
% （即真空方程）。该式说明：在矩多面体 P_\zeta 的内点上，|\phi_i|^2 等于仿射函数
% \langle\mu_F,v^i\rangle 与常数 \sum_a\zeta^a\xi_a^i 之差；令 |\phi_i|^2\ge0 即得
% \langle\mu_F,v^i\rangle\ge\sum_a\zeta^a\xi_a^i，这正是紧接着（印刷 p.37 开头）"v_i 是多面体的
% 内法向"一句的依据，(4.55) 由此而来。]
